\subsection{幂零李群}

命题 12. 设 G 为有限维李群。L(G) 为幂零的充要条件是 G 具有一个幂零开子群。

设 \(\mathbf{G}\) 具有一个幂零开子群 \(\mathbf{G}_0\)。由命题 4 和命题 6 的第 2 号推论，\(\mathcal{C}^i\mathrm{L}(G_0) = \{0\}\)（当 \(i\) 足够大时）。因此 \(\mathrm{L}(G_0) = \mathrm{L}(G)\) 为幂零的。

设 \(\mathrm{L}\langle \mathrm{G}\rangle\) 为幂零的。若 \(\mathbf{K} = \mathbf{R}\) 或 \(\mathbf{C}\)，则由命题 4 的第 2 号推论，\(\mathbf{G}_0\) 是 \(\mathbf{G}\) 的单位元分支并且是幂零的，且 \(\mathbf{G}_0\) 在 \(\mathbf{G}\) 中是开的。若 \(\mathbf{K}\) 为超度量的，则命题 6 的第 2 号推论表明，存在一个开子群 \(\mathbf{G}_1\)，它是 \(\mathbf{G}\) 的子群，还存在一个整数 \(i > 0\) 和一个邻域 \(\mathbf{V}\)，其中包含 \(e\)，且位于 \(\mathbf{G}\) 中，使得 \(\mathrm{C}^i\mathrm{G}_1\cap \mathbf{V} = \{\varepsilon \}\)。于是，若 \(\mathbf{G}_0\) 是 \(\mathbf{G}_1\) 的充分小的子群，则 \(\mathrm{C}^i\mathrm{G}_0\subset \mathrm{V}\)，从而 \(\mathrm{C}^i\mathrm{G}_0 = \{\varepsilon \}\)，且 \(\mathbf{G}_0\) 为幂零的。

设 \(\mathfrak{g}\) 为幂零李代数。对应于 \(\mathfrak{g}\) 的 Hausdorff 级数 H(X, Y) 只有有限个非零项；我们知道（第二章，第 6 节，第 5 号，注 3），合成律 \((x, y) \mapsto H(x, y)\) 在 \(\mathfrak{g}\) 上定义了一个群结构。再设 \(\mathfrak{g}\) 完备且可赋范。显然，律 H 是连续多项式（《可微与解析流形》，R，附录）。因此，赋予律 H 的 \(\mathfrak{g}\) 是一个李群 G，称为与 \(\mathfrak{g}\) 相伴的李群。由第 4 节第 2 号的引理 2 和 3，\(\mathrm{L}(\mathrm{G}) = \mathfrak{g}\)。\(\phi\) 是从 \(\mathfrak{g}\) 到 \(\mathrm{G}\) 的恒等映射，是 \(\mathrm{G}\) 的一个指数映射，满足 \(\phi(\lambda x)\phi(\lambda' x) = \phi((\lambda + \lambda') x)\) 对所有 \(x \in \mathfrak{g}\)、\(\lambda \in \mathrm{K}\)、\(\lambda' \in \mathrm{K}\) 成立。每个李子代数 \(\mathfrak{h}\)（它是 \(\mathfrak{g}\) 的）具有拓扑补空间，都是 \(\mathrm{H}\) 的一个李子群 \(\mathrm{G}\)，且 \(\mathrm{L}(\mathrm{H}) = \mathfrak{h}\)。

命题 13. 设 G 是 R 或 C 上的有限维单连通幂零李群。

\begin{enumerate}
\def\labelenumi{(\roman{enumi})}
\item
  \(\exp_{\mathbf{G}}\) 是与 \(\mathrm{L}(\mathrm{G})\) 相伴的李群到 \(\mathbf{G}\) 上的一个同构。
\item
  \(\mathbf{G}\) 的每个积分子群都是 \(\mathbf{G}\) 的单连通李子群。
\end{enumerate}

设 \(\mathfrak{g} = \mathrm{L}(\mathrm{G})\)，它是幂零的（命题 12）。由于具有同一李代数的两个 \(\mathbf{R}\) 或 \(\mathbf{C}\) 上的单连通李群是同构的（第 6 节第 3 号定理 3(ii)），只需在 \(\mathrm{G}\) 是与 \(\mathfrak{g}\) 相伴的群时证明本命题。于是 (i) 和 (ii) 由命题之前所述内容推出。

命题 14. 设 G 是 R 或 C 上的有限维连通李群。

\begin{enumerate}
\def\labelenumi{(\roman{enumi})}
\item
  若 G 为幂零的，则 \(\exp_{G}\) 是 étale 的满射。
\item
  若 K = C 且 \(\exp_{G}\) 是 étale 的，则 G 为幂零的。
\end{enumerate}

设 \(G'\) 为 G 的万有覆叠空间。设 \(\phi\) 为从 \(G'\) 到 G 的典范态射。则 \(\exp_G = \phi \circ \exp_{G'}\)（第 6 节第 4 号命题 10），故 (i) 由命题 13(i) 得出。

若 \(\mathbf{K} = \mathbf{C}\) 且 exp 是 étale 的，则对所有 \(x \in \mathrm{L}(\mathrm{G})\)，\(\operatorname{ad} x\) 都没有属于 \(2i\pi (\mathbf{Z} - \{0\})\) 的特征值（第 6 节第 4 号命题 12 的推论）。对 \(\lambda x\) 应用这一点，其中 \(\lambda\) 遍历 \(\mathbf{C}\)，可得 ad \(x\) 的所有特征值都为零，因而 ad \(x\) 是幂零的。因此 \(\mathrm{L}(\mathrm{G})\) 是幂零的（第一章第 4 节定理 1 的推论 1），从而 G 是幂零的（命题 12）。

命题 15. 设 \(\mathbf{G}\) 是 \(\mathbf{R}\) 或 \(\mathbf{C}\) 上的有限维连通幂零李群，\(\mathbf{A}\) 是 \(\mathbf{G}\) 的一个积分子群。则 \(Z_{\mathrm{G}}(\mathbf{A})\) 是 \(\mathbf{G}\) 中李代数为 \(\mathfrak{z}_{\mathfrak{g}}(\mathrm{L}(\mathbf{A}))\) 的连通李子群。

由第 3 号命题 9，只需证明 \(Z_{\mathrm{G}}(\mathbf{A})\) 是连通的。设 \(g \in Z_{\mathrm{G}}(\mathbf{A})\)。存在 \(x \in \mathrm{L}(\mathbf{G})\)，使得 \(g = \exp x\)（命题 14）。于是 \(\operatorname{Ad} g|\mathrm{L}(\mathbf{A}) = 1\)（第 3 号命题 9），因而 \(\operatorname{Ad} g^n |\mathrm{L}(\mathbf{A}) = 1\)，对所有 \(n \in \mathbf{Z}\) 成立，从而 \(\exp (\operatorname{ad} nx)|\mathrm{L}(\mathbf{A}) = 1\)，对所有 \(n \in \mathbf{Z}\) 成立。由于映射

\[
\lambda \mapsto \exp (\operatorname{ad} \lambda x) | \mathrm{L} (\mathrm{A})
\]

从 K 到 \(\mathcal{L}(\mathrm{L}(\mathrm{A}),\mathrm{L}(\mathrm{G}))\) 的映射是多项式，所以 \(\exp(\operatorname{ad}\lambda x)|\mathrm{L}(\mathrm{A})=1\) 对所有 \(\lambda\in K\) 成立，即 \(\exp(\lambda x)\in Z_{\mathrm{G}}(\mathrm{A})\) 对所有 \(\lambda\in K\) 成立。

命题 16. 设 G 是 R 或 C 上的有限维幂零李群，A 是 G 中异于 G 的一个积分子群。则 \(\mathrm{N}_{\mathrm{G}}(\mathrm{A})\) 是 G 中异于 A 的连通李子群。

\(N_{\mathrm{G}}(\mathsf{A}) \neq \mathsf{A}\)（《代数》第一章第 6 节命题 8 的推论 1）。由第 4 号命题 11，只需证明 \(N_{\mathrm{G}}(\mathsf{A})\) 是连通的。设 \(g \in N_{\mathrm{G}}(\mathsf{A})\)。存在 \(x \in \mathrm{L}(\mathrm{G})\)，使得 \(g = \exp x\)（命题 14）。设 \(\mathrm{E}\) 是 \(\mathcal{L}(\mathrm{L}(\mathrm{G}))\) 的向量子空间，由 \(u \in \mathcal{L}(\mathrm{L}(\mathrm{G}))\) 组成，且 \(u(\mathrm{L}(\mathsf{A})) \subset \mathrm{L}(\mathsf{A})\)。于是 \(\operatorname{Ad} g^n \in \mathbb{E}\)，从而 \(\exp (\operatorname{ad} nx) \in \mathbb{E}\) 对所有 \(n \in \mathbf{Z}\) 成立。因此 \(\exp (\operatorname{ad} \lambda x) \in \mathbb{E}\) 对所有 \(\lambda \in K\) 成立，即 \(\exp (\lambda x) \in N_{\mathrm{G}}(\mathsf{A})\) 对所有 \(\lambda \in K\) 成立。

命题 17. 设 \(\mathfrak{g}\) 是 \(\mathbf{K}\) 上的有限维幂零李代数，\((\mathfrak{g}_0, \mathfrak{g}_1, \ldots, \mathfrak{g}_n)\) 是 \(\mathfrak{g}\) 的递减理想序列，满足 \(\mathfrak{g}_0 = \mathfrak{g}, \mathfrak{g}_n = \{0\}\)，且 \([\mathfrak{g}, \mathfrak{g}_i] \subset \mathfrak{g}_{i+1}\) 对 \(0 \leqslant i < n\)。设 \(\mathfrak{a}_1, \mathfrak{a}_2, \ldots, \mathfrak{a}_p\) 是 \(\mathfrak{g}\) 的向量子空间，并且每个 \(\mathfrak{g}_i\) 都是它与各个 \(\mathfrak{a}_j\) 的交的直和。赋予 \(\mathfrak{g}\) Hausdorff 合成律 H。设 \(\phi\) 是映射

\[
\left(x _ {1}, x _ {2}, \dots , x _ {p}\right) \mapsto x _ {1} \mathsf {H} x _ {2} \mathsf {H} \dots \mathsf {H} x _ {p}
\]

\[
o f a _ {1} \times a _ {2} \times \dots \times a _ {p}
\]

\begin{enumerate}
\def\labelenumi{(\roman{enumi})}
\item
  \(\phi\) 是从 \(a_{1} \times a_{2} \times \cdots \times a_{p}\) 到 \(g\) 的双射；
\item
  \(\phi\) 和 \(\phi^{-1}\) 都是多项式映射；
\item
  映射 \((x,y)\mapsto\phi^{-1}(\phi(x)\cdot\phi(y)^{-1})\)，从 \((\mathfrak{a}_{1}\times\mathfrak{a}_{2}\times\cdots\times\mathfrak{a}_{p})^{2}\) 到 \(a_{1}\times a_{2}\times\cdots\times a_{p}\)，是多项式。
\end{enumerate}

当 \(\dim \mathfrak{g} = 0\) 时本命题显然成立。设 \(\dim \mathfrak{g} > 0\)，并且本命题已对维数 \(<\dim \mathfrak{g}\) 成立。可以假定 \(\mathfrak{g}_{n-1} \neq \{0\}\)，于是 \(\mathfrak{g}_{n-1}\) 是 \(\mathfrak{g}\) 的非零中心理想。存在指标 \(j\)，使得 \(\mathfrak{b} = \mathfrak{g}_{n-1} \cap \mathfrak{a}_j \neq \{0\}\)。令 \(\mathfrak{g}' = \mathfrak{g}/\mathfrak{b}\)，\(\theta\) 为从 \(\mathfrak{g}\) 到 \(\mathfrak{g}'\) 的典范态射，\(\mathfrak{g}'_i = \theta(\mathfrak{g}_i)\) 且 \(\mathfrak{a}_i' = \theta(\mathfrak{a}_i)\)。则 \((\mathfrak{g}_0', \mathfrak{g}_1', \ldots, \mathfrak{g}_n')\) 是 \(\mathfrak{g}'\) 的递减理想序列，满足 \(\mathfrak{g}_0' = \mathfrak{g}'\)，\(\mathfrak{g}_n' = \{0\}\)，

\[
[ \mathfrak {g} ^ {\prime}, \mathfrak {g} _ {i} ^ {\prime} ] \subset \mathfrak {g} _ {i + 1} ^ {\prime}
\]

当 \(0 \leqslant i < n\) 时，且每个 \(\mathfrak{g}_i'\) 都是它与各个 \(\mathfrak{a}_j'\) 的交的直和。设 \(\phi'\) 是映射

\[
\left(x _ {1} ^ {\prime}, x _ {2} ^ {\prime}, \dots , x _ {p} ^ {\prime}\right) \mapsto x _ {1} ^ {\prime} \mathsf {H} x _ {2} ^ {\prime} \mathsf {H} \dots \mathsf {H} x _ {p} ^ {\prime}
\]

从 \(a_{1}^{\prime} \times a_{2}^{\prime} \times \cdots \times a_{p}^{\prime}\) 到 \(g^{\prime}\) 的映射。由归纳假设，\(\phi^{\prime}\) 是双射，且 \(\phi^{\prime}, \phi^{\prime-1}\) 是多项式映射。

设 \(x \in g\)。记

\[
\phi^ {\prime - 1} (\theta (x)) = \left(x _ {1} ^ {\prime} (x), x _ {2} ^ {\prime} (x), \dots , x _ {p} ^ {\prime} (x)\right).\tag{1}
\]

于是

\[
\theta (x) = x _ {1} ^ {\prime} (x) \mathsf {H} x _ {2} ^ {\prime} (x) \mathsf {H} \dots \mathsf {H} x _ {p} ^ {\prime} (x).\tag{2}
\]

设 \(\mathfrak{h}_1\) 是 \(\mathfrak{b}\) 在 \(\mathfrak{g}\) 中的一个向量补空间，它等于各个 \(a_k\)（当 \(k \neq j\) 时）之和与 \(\mathfrak{b}\) 在 \(a_j\) 中的一个补空间之和。存在从 \(\eta\) 到 \(\mathfrak{g}'\) 的双射 \(\mathfrak{h}_1\)，满足 \(\theta \circ \eta = \mathrm{Id}_{\mathfrak{g}'}\)。对所有 \(x \in \mathfrak{g}\)，记

\[
\zeta (x) = \eta (x _ {1} ^ {\prime} (x)) \mathsf {H} \eta (x _ {2} ^ {\prime} (x)) \mathsf {H} \dots \mathsf {H} \eta (x _ {p} ^ {\prime} (x)) \in \mathfrak {g}.\tag{3}
\]

\[
y (x) = \zeta (x) ^ {- 1} \text { H } x = (- \zeta (x))\cdot x.\tag{4}
\]

由 (2) 和 (3)，\(\theta (\zeta (x)) = \theta (x)\)，从而 \(y(x)\in \mathfrak{h}\)。最后记

\[
\psi (x) = \left(\eta \left(x _ {1} ^ {\prime} (x)\right), \dots , \eta \left(x _ {j} ^ {\prime} (x)\right) + y (x), \dots , \eta \left(x _ {p} ^ {\prime} (x)\right)\right) \in a _ {1} \times \dots \times a _ {p}.\tag{5}
\]

由于 \(y(x)\) 在 \(\mathfrak{g}\) 的中心中，

\[
\begin{array}{l l} \phi (\psi (x)) = \eta (x _ {1} ^ {\prime} (x)) \texttt {H} \dots \texttt {H} \eta (x _ {j} ^ {\prime} (x)) \texttt {H} \dots \texttt {H} \eta (x _ {p} ^ {\prime} (x)) \texttt {H} y (x) \\ = \zeta (x) \texttt {H} y (x) & \text {by (3)} \\ = x & \text {by (4)}. \end{array}
\]

因此 \(\phi \circ \psi = \mathrm{Id}_{\mathfrak{g}'}\)。现在设 \((x_{1}, x_{2}, \ldots, x_{p}) \in \mathfrak{a}_{1} \times \mathfrak{a}_{2} \times \dots \times \mathfrak{a}_{p}\)，并记 \(x = \phi(x_{1}, x_{2}, \ldots, x_{p}) = x_{1} \mathsf{H} x_{2} \mathsf{H} \dots \mathsf{H} x_{p}\)。于是

\[
\theta (x) = \theta (x _ {1}) \mathsf {H} \theta (x _ {2}) \mathsf {H} \dots \mathsf {H} \theta (x _ {p}),
\]

所以 \(x_{i}^{\prime}(x) = \theta (x_{i})\) 对所有 \(1\leqslant i\leqslant p\) 成立，从而

\[
\begin{array}{c} \zeta (x) = x _ {1} \mathsf {H} x _ {2} \mathsf {H} \dots \mathsf {H} (\eta \theta (x _ {j})) \mathsf {H} \dots \mathsf {H} x _ {p} \\ y (x) = x _ {j} - \eta \theta (x _ {j}). \end{array}
\]

由 (5)，

\[
\psi (x) = \left(x _ {1}, \dots , \eta \theta \left(x _ {j}\right) + x _ {j} - \eta \theta \left(x _ {j}\right), \dots , x _ {p}\right) = \left(x _ {1}, x _ {2}, \dots , x _ {p}\right).
\]

因此 \(\psi \circ \phi = \mathrm{Id}_{\mathfrak{a}_1 \times \ldots \times \mathfrak{a}_p}\)。这就证明了 (i)。由于 Hausdorff 律是多项式的，\(\phi\) 是多项式映射。由归纳假设，\(\phi'^{-1}\) 是多项式映射；根据公式 (1)，函数 \(x_j'\) 是多项式的，因而 \(\zeta\) 是多项式映射（公式 (3)），\(y\) 是多项式的（公式 (4)），而 \(\psi\) 是多项式映射（公式 (5)）。这就证明了 (ii)。(iii) 由 (i)、(ii) 以及 Hausdorff 律为多项式这一事实推出。

幂零李群的例子。设 G 是 \(\mathbf{GL}(n,\mathbf{K})\) 的严格下三角子群。它是 \(\mathbf{GL}(n,\mathbf{K})\) 的李子群，且 \(\mathrm{L}(\mathrm{G})\subset\mathrm{gl}(n,\mathbf{K})\) 是对角线为零的下三角矩阵李代数（第 3 节第 10 号命题 36）。由第二章第 4 节第 6 号注，G 是幂零的。以下设 K=R 或 C。由于 G 同胚于 \(\mathbf{K}^{n(n-1)/2}\)，G 是单连通的。L(G) 到 G 的指数映射就是映射

\[
u \mapsto \exp u = \sum_ {k \geqslant 0} \frac {u ^ {k}}{k !} = \sum_ {k = 0} ^ {n - 1} \frac {u ^ {k}}{k !}
\]

（第 6 节第 4 号例）。由命题 13，指数映射是从流形 L(G) 到流形 G 的同构。第 6 节第 9 号命题 17 给出其逆双射 log。在 \(\mathbf{K}^{n}\) 上赋予一个范数。由《谱理论》第一章第 4 节第 9 号，当 g ∈ G 且 \textbar\textbar g − 1\textbar\textbar{} \textless{} 1 时，

\[
\log g = \sum_ {k \geqslant 1} \frac {(- 1) ^ {k - 1}}{k} (g - 1) ^ {k}
\]

即

\[
\log g = \sum_ {k = 1} ^ {n - 1} \frac {(- 1) ^ {k - 1}}{k} (g - 1) ^ {k}.\tag{6}
\]

但 (6) 的两边都是关于 \(g\) 的解析函数，对 \(g \in \mathbf{G}\) 成立，因而对所有 \(g \in \mathbf{G}\) 都相等。

命题 18. 设 \(k\) 是交换域。V 是维数 \(>0\) 的 \(k\) 上向量空间，且维数大于 0，G 是 \(\mathbf{GL}(\mathbf{V})\) 的一个子群，其元素都是幺幂的。

\begin{enumerate}
\def\labelenumi{(\roman{enumi})}
\item
  存在一个非零元素 \(v\)，属于 \(V\)，使得 \(gv = v\) 对所有 \(g \in G\) 成立。
\item
  存在 V 的一个基 B，使得对所有 \(g \in G\)，\(g\) 关于 B 的矩阵都是下三角的，且其所有对角元都等于 1。
\item
  群 G 是幂零的。
\end{enumerate}

\begin{enumerate}
\def\labelenumi{(\alph{enumi})}
\tightlist
\item
  首先设 \(k\) 是代数闭域，且 \(G\) 的恒等表示是单的。设 \(a, b\) 属于 \(G\)。于是
\end{enumerate}

\[
\operatorname{Tr} (a (b - 1)) = \operatorname{Tr} (a b - 1) - \operatorname{Tr} (a - 1) = 0 - 0 = 0
\]

因为 \(ab - 1\) 和 \(a - 1\) 都是幂零的。由于由 \(\mathcal{L}(\mathrm{V})\) 中的 \(\mathbf{G}\) 生成的向量子空间就是 \(\mathcal{L}(\mathrm{V})\)（《代数》第八章第 4 节命题 2 的推论 1），\(\operatorname{Tr}(u(b - 1)) = 0\) 对所有 \(u \in \mathcal{L}(\mathrm{V})\) 成立，从而 \(b = 1\)。因此 \(\mathrm{G} = \{1\}\)。

\begin{enumerate}
\def\labelenumi{(\alph{enumi})}
\setcounter{enumi}{1}
\item
  现在转到一般情形。设 \(\overline{k}\) 是 \(k\) 的代数闭包，\(\overline{\mathrm{V}} = \mathrm{V} \otimes_k \overline{k}\)，且 \(\overline{\mathbf{G}} \subset \mathbf{GL}(\overline{\mathrm{V}})\) 是由 \(a \otimes 1\)（\(a \in \mathbf{G}\)）组成的集合。令 \(\mathrm{W}\)（分别地，\(\mathrm{W}'\)）为 \(\mathrm{V}\)（分别地，\(\overline{\mathrm{V}}\)）中在 \(\mathrm{G}\)（分别地，\(\overline{\mathbf{G}}\)）下不变的元素的集合。则 \(\mathrm{W}' = \mathrm{W} \otimes_k \overline{k}\)，其中 \(\mathrm{W} = \bigcap_{g \in \mathbb{G}} \mathrm{Ker}(g - 1)\)，并且 \(\mathrm{W}' = \bigcap_{g \in \mathbb{G}} \mathrm{Ker}(g - 1) \otimes 1\)。若 \(\mathrm{V}_1\) 表示 \(\overline{\mathrm{V}}\) 的在 \(\overline{\mathbf{G}}\) 下稳定的非零向量子空间集合中的极小元，则由证明的 (a) 部分，\(\mathrm{V}_1 \subset \mathrm{W}'\)；因此 \(\mathrm{W} \neq \{0\}\)，这就证明了 (i)。
\item
  对 dim V 作归纳，由 (i) 可得，存在 V 的一个在 G 下稳定的向量子空间递增序列 \(\left(\mathrm{V}_{1},\mathrm{V}_{2},\ldots,\mathrm{V}_{n}\right)\)，使得 \(V_{n}=V\)，并且由 G 典范地导出的 \(V_{i}/V_{i-1}\) 的自同构群对所有 i 都退化为 \(\{1\}\)（约定 \(V_{r}=\{0\}\) 对 \(r\leqslant0\) 成立）。这首先推出 (ii)，继而推出 (iii)（第二章第 4 节第 6 号注）。
\end{enumerate}

推论 1. 设 G 是有限维连通实或复李群。G 为幂零的充要条件是 Ad G 的每个元素都是幂零的。

若 Ad G 的每个元素都是幺幂的，则 Ad G 是幂零的（命题 18），因而作为 Ad G 的中心扩张的 G 是幂零的。若 G 是幂零的，则 L(G) 是幂零的，所以对所有 \(x \in \mathrm{L}(G)\)，ad x 都是幂零的，从而 \(\mathrm{Ad}(\exp x) = \exp \mathrm{ad} x\) 是幺幂的；但 G 的每个元素都具有 \(\exp x\) 的形式，其中 \(x \in \mathrm{L}(G)\)（命题 14）。

推论 2. \(\mathbf{GL}(n, \mathbf{K})\) 的每个由幺幂元素组成的解析子群都是单连通李子群。

这由命题 13(ii)、命题 18(ii) 以及严格下三角群是单连通的这一事实推出。

